\documentclass[dvipdfmx]{jarticle}
\pagestyle{empty}
\hoffset=-25mm \voffset=-25mm
\setlength{\textwidth}{16cm}
\setlength{\textheight}{24cm}
\usepackage{url}
\usepackage{tikz}
\usetikzlibrary{calc,shapes}
\usepackage{fancyvrb}
\usepackage{tabularx}
\usepackage{multirow}
\newcounter{coffee}
\newcounter{bean}
\newcommand{\ttdo}{\texttt{do}}
\newcommand{\ttwhile}{\texttt{while}}
\newcommand{\ttlp}{\texttt{(}}
\newcommand{\ttrp}{\texttt{)}}
\newcommand{\ttsc}{\texttt{;}}
\newcommand{\tteq}{\texttt{=}}
\newcommand{\ttif}{\texttt{if}}
\newcommand{\ttelse}{\texttt{else}}
\newcommand{\ttand}{\texttt{\&\&}}
\newcommand{\ttlt}{\texttt{<}}
\newcommand{\ttsub}{\texttt{-}}

\newcommand{\newLabel}[1]{\texttt{newLabel}({#1})}
\newcommand{\newTemp}[1]{\texttt{newTemp}({#1})}

\newcommand{\redtext}[1]{{\color{red}{#1}}}
\begin{document}
\begin{flushright}
\verb|2026/11/6|
\end{flushright}
\begin{center}{\bf\Large
  2026年度 コンパイラ演習課題(6)
}\end{center}
\begin{itemize}
  \item TACTの課題の添付ファイルとしてPDFファイルで提出すること．
  \item 手書きレポートをPDFとして提出する方法は講義webサイトを
  参照\footnote{\url{https://sites.google.com/sqlab.jp/26compiler/how-to-make-report}}すること．
  \item 提出期限: {11月13日（金）23:55}
\end{itemize}

\vspace*{1cm}

\begin{list}{{\bf [\,\arabic{coffee}\,]}}{\usecounter{coffee}}
\item
以下の属性文法$G_1$の規則を考える．ここで$N,L,B$は非終端記号，$0,1,\texttt{.}$は終端記号である．
$N$が開始記号である．
(これは２進数固定小数点の値を与える)

\[
\begin{array}{lll}
  N\rightarrow L^{(1)}\;\texttt{.}\;L^{(2)} & \{N.val=L^{(1)}.val+L^{(2)}.val, & L^{(1)}.s=0,\\
  & & L^{(2)}.s=-L^{(2)}.d\}\\
  L\rightarrow L^{(1)}\;B & \{L.val = L^{(1)}.val+B.val, & L^{(1)}.s=L.s+1,\\
    & \phantom{\{}L.d=L^{(1)}.d+1, & B.s=L.s\}\\
  L\rightarrow B & \{L.val=B.val, & B.s=L.s,\\
    & \phantom{\{}L.d=1\} & \\
  B\rightarrow 0 & \{B.val=0\}\\
  B\rightarrow 1 & \{B.val=2^{B.s}\}
\end{array}\]

\begin{list}{(\alph{bean})}{\usecounter{bean}}
  \item $X.val$, $X.s$, $X.d$のうち，どれが合成属性であるか示せ．ここで，$X\in\{N,L,B\}$とする．
  \item \texttt{101.011}に対する$N.val$の値を求めよ．
\end{list}
\item 以下の文法$G_2$の生成規則と意味規則を考える．
\[
\begin{array}{ll}
  E\rightarrow TE' & \{E'.a=T.val,\ E.val=E'.b\}\\
  E'\rightarrow +TE'^{(1)} & \{{E'^{(1)}.a=E'.a+T.val},\ E'.b=E'^{(1)}.b\}\\
  E'\rightarrow\varepsilon & \{E'.b=E'.a\}\\
  T\rightarrow FT' & \{T'.a=F.val,\ T.val=T'.b\}\\
  T'\rightarrow \ast FT'^{(1)} & \{T'^{(1)}.a=T'.a\times F.val,\ T'.b=T'^{(1)}.b\}\\
  T'\rightarrow\varepsilon & \{T'.b=T'.a\}\\
  F\rightarrow (E) & \{F.val=E.val\}\\
  F\rightarrow\mathsf{NUM} & \{F.val=\mathsf{NUM}.val\}
\end{array}
\]
  \begin{list}{(\alph{bean})}{\usecounter{bean}}
    \item $X.val$, $X.a$, $X.b$のうち，どれが合成属性であるか示せ．ここで，$X\in\{E,E',T,T',F\}$とする．
    \item $4+6\ast 7$に対する$E.val$の値を上記の意味規則によって評価せよ．
    \item $(4+6)\ast 7$に対する$E.val$の値を上記の意味規則によって評価せよ．
  \end{list}
\end{list}
\end{document}